Amortized inference in cognitive modeling requires an explicit statement of what an estimator is amortized over.
We organize models into a three-level taxonomy---\textbf{model species}, \textbf{model family}, and \textbf{model class} (see also \autoref{fig:levels})---in which each level comprises multiple members of the level below.
The taxonomy fixes the level at which inference is amortized and, with it, the generalization scope of a trained estimator.

A \textbf{model species} $\mathcal{M}$ has a fixed computational structure with all parameters fully specified.
Standard SBI in cognitive modeling amortizes over a single species: any change in parameterization or generative structure requires retraining.

A \textbf{model family} $\mathcal{F}$ comprises species that share a computational structure but differ in parameterization or auxiliary assumptions.
Variants of the classic diffusion decision model \citep[DDM;][]{ratcliff1978diffusionModel, ratcliffDiffusionModel2004} differ in collapsing versus fixed boundaries \citep{bogaczPhysicsOptimal2006, ditterichEvidenceTimevariant2006, rasanan2025collapsing}, starting-point variability, non-decision time structure, or which parameters carry regressors.
Such specifications are typically nested---simpler variants follow from more general ones by fixing parameters to zero or another constant---so a single network can in principle cover an entire family.

A \textbf{model class} $\mathcal{C}$ groups families that implement a shared computational principle.
In decision making, the DDM, the racing diffusion model \citep[RDM;][]{tillman2020racingDiffusion, zandbeltResponseTimes2014}, and the circular diffusion model \citep[CDM;][]{smith2016circularDiffusion, rasananDiscretechoiceOptions2024} together form the class of sequential sampling models (SSMs).
These families are not nested and cannot be reached from one another by parameter constraints.
Class-level amortization therefore demands generalization across structurally distinct generative mechanisms while exploiting what they share, such as common target observables.

We demonstrate family- and class-level amortization under the generative meta-model
%
\begin{equation}\label{eq:generative}
\mathcal{F} \sim p(\mathcal{F} \mid \mathcal{C}), \,\,
(\mathcal{M},\,\,\mathcal{D}) \sim p(\mathcal{M},\,\, \mathcal{D} \mid \mathcal{F}), \,\,
\Theta_\mathcal{M} \sim p(\Theta_\mathcal{M} \mid \mathcal{M}, \mathcal{D}), \,\,
\rmY \sim p(\rmY \mid \Theta_\mathcal{M},\mathcal{M}, \mathcal{D}),
\end{equation}
%
\noindent where $\mathcal{D}$ is a design configuration over the parameters $\Theta_\mathcal{M}$ of species $\mathcal{M}$ and $\rmY$ is observed behavior.
The sections that follow describe how the model and design indices $\mathcal{M}$ and $\mathcal{D}$ are represented numerically, through embeddings of positional encodings and design matrices, and how the CogFormer architecture maps them to the amortized posterior $p(\Theta_{\mathcal{M}} \mid \rmY, \mathcal{M}, \mathcal{D})$.
